\subsection{Artifact Coverage and Reuse}
\label{app:artifact-scope}
The released SRA-Bench task package is documented as English. This metadata does not imply an English-only retrieval corpus: web-collected skill text includes non-Latin scripts, and the corpus has no per-record language labels. The six benchmarks cover theorem application, logical reasoning, tool workflows, medical calculation, mathematical problem solving, and code generation. Our evaluation does not establish multilingual performance or demographic representativeness.

The MIT declaration in the SRA-Bench metadata does not supersede the rights of constituent artifacts. MedCalc-Bench data carry a CC-BY-SA-4.0 notice, while individual web-collected skills include CC-BY-SA, CC-BY-NC-SA, and GPL notices. Code and content permissions can also differ: a benchmark's software license does not necessarily cover textbook questions, source corpora, or every inherited evaluator. Reuse and redistribution must respect these component-specific conditions rather than treating the complete collection as uniformly MIT-licensed.

The task package includes medical-calculation questions, and the accompanying ToolQA source assets include host and location fields. Public availability does not certify that these materials are free of identifying information. Credential-file exclusion from a software archive is distinct from reviewing and de-identifying task content or model traces. Similarly, routing labels and held-out validation measure skill-selection behavior, not objectionable-content screening. The retrieved instructions and executable examples remain untrusted content; routing performance does not certify their safety.
